\section{Giris}

Gercek zamanli konum ve durum izleme sistemleri, modern toplu tasima altyapilarinin en gorunur bilesenlerinden biridir. Yolcularin otobusun mevcut konumunu, tahmini varis suresini ve hatta yogunluk seviyesini takip edebilmesi, hizmet kalitesini dogrudan etkileyen bir unsurdur. Benzer sekilde, isletmeci kurumlar acisindan da canli telemetri verisi; hat optimizasyonu, gecikme tespiti ve operasyonel izleme icin yuksek deger tasir. Bu nedenle toplu tasimada gercek zamanli veri akisi, yalnizca bir bilgi ekrani problemi degil, ayni zamanda bulut tabanli veri isleme ve servis butunlestirme problemidir.

Bu projede, gercek EGO verisi veya fiziksel cihazlar kullanilmadan, EGO-benzeri bir otobus takip sistemi sentetik veri ile modellenmistir. Buradaki temel amac, gercek bir kurum verisine bagli kalmadan, otobus telemetrisinin MQTT tabanli bir yapi ile buluta aktarilmasi, bulutta islenmesi ve bir dashboard uzerinden canli olarak izlenmesi surecini gostermektir. Bu tercih, dersin odak noktasini veri kaynagindan cok mimari tasarim, servis entegrasyonu ve gercek zamanli akis yonetimi uzerine tasimaktadir. Ayrica sentetik veri kullanimi, sistem tasarimini test etmeyi kolaylastirirken gercek kurum verisine bagimlilik gibi operasyonel kisitlari da ortadan kaldirmaktadir.

Calismanin ana katkisi, veri uretimi, bulut uzerinde zenginlestirme ve canli sunum katmanlarini tek bir yapi altinda birlestirmesidir. Simulator tarafinda yalnizca ham telemetri verisi uretilmis; ETA, tahmini inis miktari, doluluk ve gecikme gibi karar verici bilgiler ise AWS uzerinde hesaplanmistir. Boylece sistem, sadece sabit veri gosteren bir demo olmaktan cikmis, gercek zamanli olay akisina dayali bulut tabanli bir referans mimariye donusmustur. MQTT'nin hafif mesajlasma yapisi \cite{mqtt311}, AWS IoT Core'un cihaz-bulut iletisim destegi \cite{aws_iot_mqtt} ve Kinesis Data Streams'in gercek zamanli veri yutma kapasitesi \cite{kinesis_intro} bu tasarimin temelini olusturmaktadir.
